%------------------------------------------------------------------------------%
% Luis A. D'Afonseca
%
%------------------------------------------------------------------------------%
\documentclass[a4paper,12pt,fleqn]{article}

\usepackage{style_avaliacao}

\ShowAnswers

\newcommand{\D}[2]{\frac{\partial {#1}}{\partial {#2}}}
\newcommand{\DS}[2]{\frac{\partial^2 {#1}}{\partial {#2}^2}}
\newcommand{\DM}[3]{\frac{\partial^2 {#1}}{\partial {#2} \partial {#3}}}


\newcommand{\vetor}[2]{\left(\begin{array}{c} #1 \\ #2 \end{array}\right)}

\newcommand{\veci}{\bm{i}}
\newcommand{\vecj}{\bm{j}}
\newcommand{\veck}{\bm{k}}

%------------------------------------------------------------------------------%
\begin{document}

\makeHeader{CFVVI}{Exame Especial -- Turma 4}{18/09/2024}

% 01
%------------------------------------------------------------------------------%
\exercise{25}
Calcule o limite solicitado, ou prove que o limite não existe
\(
  \lim_{(x, y)\to(1, 1)} \frac{x^2 - y^2}{y-x}
\)

\begin{answer}
  Se tentarmos calcular diretamente a fração obtemos uma
  indeterminação $\sfrac{0}{0}$.
  Portanto, precisamos remover a indeterminação manipulando
  algebricamente a função,
  assumindo que $(x, y)\neq(1, 1)$, temos
  \begin{align*}
    f(x) & = \frac{x^2 - y^2}{y-x}   \\[2mm]
         & = \frac{(x-y)(x+y)}{y-x}  \\[2mm]
         & = -\frac{(x-y)(x+y)}{x-y} \\[2mm]
         & = -(x+y)                  \\[2mm]
         & = -x-y
  \end{align*}
  Onde o cancelamento só foi possível, pois $(x, y)\neq(1, 1)$.
  Agora podemos calcular o limite, pois temos uma função contínua em $(1, 1)$,
  \[
    L
    = \lim_{(x, y)\to(1, 1)} \frac{x^2 - y^2}{y-x}
    = \lim_{(x, y)\to(1, 1)} -x-y = -2
  \]
  \clearpage
\end{answer}

% 02
%------------------------------------------------------------------------------%
\exercise{25}
Encontre linearização da função
\(
  f(x, y) = x^2 - xy - y^2
\)
no ponto \((1, 1)\).

\begin{answer}
  A linearização de $f$ no ponto \((1, 1)\) é a função
  \begin{align*}
    L(x, y)
    & = f(x_0, y_0) + f_x(x_0, y_0)(x-x_0) + f_y(x_0, y_0)(y-y_0) \\
    & = f(1, 1) + f_x(1, 1)(x-1) + f_y(1, 1)(y-1)
  \end{align*}
  Precisamos das derivadas parciais de $f$
  \[
    f_x(x, y)
    = \D{f}{x}
    = \D{}{x}\left[x^2 - xy - y^2\right]
    = 2x -y
  \]
  \[
    f_y(x, y)
    = \D{f}{y}
    = \D{}{y}\left[x^2 - xy - y^2\right]
    = -x -2y
  \]
  Avaliando $f$ e suas derivadas no ponto $(1, 1)$
  \[
    f(1, 1)
    = \left(x^2 - xy - y^2\right)\evalat{(1, 1)}{}
    = 1^2 - 1\times 1 - 1^2
    = -1
  \]
  \[
    f_x(1, 1)
    = \left(2x -y\right)\evalat{(1, 1)}{}
    = 2\times 1 - 1
    = 1
  \]
  \[
    f_y(1, 1)
    = \left(-x -2y\right)\evalat{(1, 1)}{}
    = -1 -2 \times 1
    = -3
  \]
  Assim
  \begin{align*}
    L(x, y)
    & = f(1, 1) + f_x(1, 1)(x-1) + f_y(1, 1)(y-1) \\
    & = -1 + 1(x-1) -3(y-1) \\
    & = -1 + x -1 -3y +3 \\
    & = x - 3y + 1
  \end{align*}
  \clearpage
\end{answer}

% 03 - pg 270 ex 14 - res pg 837
%------------------------------------------------------------------------------%
\exercise{25}
Encontre e classifique os pontos críticos da função
\(
  f(x, y) = x^3 + 3xy + y^3
\)

\begin{answer}

  Precisamos das derivadas parciais de $f$
  \[
    \D{f}{x}
    = \D{}{x}\left[x^3 + 3xy + y^3\right]
    = 3x^2 + 3y
  \]
  \[
    \D{f}{y}
    = \D{}{y}\left[x^3 + 3xy + y^3\right]
    = 3x + 3y^2
  \]
  então
  \[
    \nabla f = \vetor{3x^2 + 3y}{3x + 3y^2}
  \]

  A função é um polinômio, então possui derivadas em todos os pontos do plano.
  Assim os pontos críticos são apenas os pontos onde as derivadas parciais são zero,
  \(\nabla f = 0\)
  \[
    3x^2 + 3y = 0
    \qquad\text{e}\qquad
    3x + 3y^2 = 0
  \]
  ou, simplificando,
  \[
    x^2 + y = 0
    \qquad\text{e}\qquad
    x + y^2 = 0
  \]
  Isolando $y$ na primeira equação, \(y = -x^2\), e substituindo na segunda, temos
  \begin{align*}
    x + y^2                 & = 0 \\
    x + \left(-x^2\right)^2 & = 0 \\
    x + x^4                 & = 0 \\
    x(1+ x^3)               & = 0
  \end{align*}
  As soluções dessa equação são $x=0$ ou $x=-1$.
  Se $x=0$ temos $y=0$ e
  se $x=-1$ temos $y=-1$.
  Portanto, os pontos críticos são
  \[
    (x_1, y_1) = (0, 0)
    \qquad\text{e}\qquad
    (x_2, y_2) = (-1, -1)
  \]

  Precisamos das derivadas parciais de segunda ordem de $f$
  \[
    \DS{f}{x}
    = \D{}{x}\left[3x^2 + 3y\right]
    = 6x
  \]
  \[
    \D{f}{y}
    = \D{}{y}\left[3x + 3y^2\right]
    = 6y
  \]
  \[
    \DM{f}{x}{y}
    = \D{}{x}\D{f}{y}
    = \D{}{x}\left[3x + 3y^2\right]
    = 3
  \]
  então
  \[
    H = \left(\begin{array}{cc}
      6x & 3 \\
      3  & 6y
    \end{array}\right)
  \]

  Para classificar os pontos críticos precisamos avaliar o discriminante nos
  pontos críticos.

  Considerando o ponto \((x_1, y_1) = (0, 0)\)
  \[
    f_{xx}(0, 0) = 0
  \]
  \[
    f_{yy}(0, 0) = 0
  \]
  \[
    f_{xy}(0, 0) = 3
  \]
  \[
    D_1
    = f_{xx}(0, 0)f_{yy}(0, 0) - f^2_{xy}(0, 0)
    = 0 \times 0 - 3^2
    = -9 < 0
  \]
  Portanto, o ponto $(0, 0)$ é um ponto de sela.

  Considerando o ponto \((x_2, y_2) = (-1, -1)\)
  \[
    f_{xx}(-1, -1) = -6
  \]
  \[
    f_{yy}(-1, -1) = -6
  \]
  \[
    f_{xy}(-1, -1) = 3
  \]
  \[
    D_2
    = f_{xx}(-1, -1)f_{yy}(-1, -1) - f^2_{xy}(-1, -1)
    = (-6)(-6) - 3^2
    = 36 - 9
    = 25 > 0
  \]
  Portanto, o ponto $(-1, -1)$ é um máximo ou mínimo local.
  Como \(f_{xx}(-1, -1) = -6 < 0\)
  o ponto é um ponto de máximo local.
  \clearpage
\end{answer}


% 04
%------------------------------------------------------------------------------%
\exercise{25}
Encontre os valores máximo e mínimo de \(f(x, y, z) = x - 2y + 5z\)
na esfera \(x^2+y^2+z^2=30\)

\begin{answer}
  Queremos os extremos da função
  \[
    f(x, y, z) = x - 2y + 5z
  \]
  restrita a
  \[
    g(x, y, z) = x^2 + y^2 + z^2 = 30
  \]
  Para aplicarmos multiplicadores de Lagrange precisamos das derivadas parciais
  das funções $f$ e $g$
  \[
    f_x(x, y, z) =  1 \qquad\qquad
    f_y(x, y, z) = -2 \qquad\qquad
    f_z(x, y, z) =  5
  \]
  \[
    g_x(x, y, z) = 2x \qquad\qquad
    g_y(x, y, z) = 2y \qquad\qquad
    g_z(x, y, z) = 2z
  \]
  Precisamos resolver o sistema \(\nabla f = \lambda \nabla g\) e \(g=30\)
  \[
  \begin{cases}
    \hphantom{-}1 = 2\lambda x \\
    -2 = 2\lambda y \\
    \hphantom{-}5 = 2\lambda z \\
    x^2 + y^2 + z^2 = 30
  \end{cases}
  \]
  Da primeiro equação temos
  \[
    \lambda = \frac{1}{2x}
  \]
  Isolando $y$ na segunda e substituindo $\lambda$ temos
  \[
    y
    = \frac{-2}{2\lambda}
    = \frac{-1}{\lambda}
    = -2x
  \]
  Fazendo o mesmo na terceira temos
  \[
    z
    = \frac{5}{2\lambda}
    = \frac{5\times 2x}{2}
    = 5x
  \]
  Substituindo $y$ e $z$ na quarta equação
  \begin{align*}
    x^2 + y^2 + z^2        & = 30 \\
    x^2 + (-2x)^2 + (5x)^2 & = 30 \\
    x^2 + 4x^2 + 25x^2     & = 30 \\
    30 x^2 = 30            &  \\
    x = \pm 1              &
  \end{align*}

  Para $x=1$
  \[
    y = -2x = -2 \qquad\qquad
    z = 5x = 5
  \]
  portanto
  \[
    f(1, -2, 5)
    = 1 - 2(-2) + 5\times 5
    = 1 + 4 + 25
    = 30
  \]

  Para $x=-1$
  \[
    y = -2x = 2 \qquad\qquad
    z = 5x = -5
  \]
  portanto
  \[
    f(-1, 2, -5)
    = -1 - 2\times 2 + 5(-5)
    = -1 - 4 - 25
    = -30
  \]

  O valor mínimo de $f$ é $-30$ e o valor máximo é $30$
\end{answer}

%------------------------------------------------------------------------------%
\end{document}
%------------------------------------------------------------------------------%
